\subsection*{Procedimiento y operaciones del simulacro S08}
Se invierte la ecuación de Basquin, se suman fracciones de vida (Miner), se aplican las reglas de mezclas en paralelo y en serie y se usa $K_I=Y\sigma\sqrt{\pi a}$ igualada a $K_{IC}$.

\begin{enumerate}

\item Coeficiente de resistencia a fatiga.

\[\fitmath{\mathrm{sf}=900\approx 900\,\mathrm{MPa}}\]

\item Exponente de Basquin.

\[\fitmath{b=-0{,}09\approx -0{,}09\,}\]

\item Amplitud de tensión.

\[\fitmath{\mathrm{sa}=300\approx 300\,\mathrm{MPa}}\]

\item Reversiones hasta rotura.

\[\fitmath{N_{2}={\left(\frac{\mathrm{sa}}{\mathrm{sf}}\right)}^{\frac{1}{b}}}\]

\[\fitmath{N_{2}={\left(\frac{(300)}{(900)}\right)}^{\frac{1}{(-0{,}09)}}\approx 200146{,}2\,}\]

\item Ciclos hasta rotura.

\[\fitmath{\mathrm{Nf}=\frac{N_{2}}{2}}\]

\[\fitmath{\mathrm{Nf}=\frac{(200146{,}1658)}{2}\approx 100073{,}1\,\mathrm{ciclos}}\]

\item Daño del bloque a la amplitud mayor.

\[\fitmath{d_{1}=\frac{20000}{100000}\approx 0{,}2\,}\]

\item Daño del bloque a la amplitud intermedia.

\[\fitmath{d_{2}=\frac{50000}{500000}\approx 0{,}1\,}\]

\item Daño del bloque a la amplitud menor.

\[\fitmath{d_{3}=\frac{100000}{2000000}\approx 0{,}05\,}\]

\item Daño de Miner por bloque.

\[\fitmath{\mathrm{Dm}=d_{1}+d_{2}+d_{3}}\]

\[\fitmath{\mathrm{Dm}=(0{,}2)+(0{,}1)+(0{,}05)\approx 0{,}35\,}\]

\item Bloques hasta D=1.

\[\fitmath{\mathrm{nb}=\frac{1}{\mathrm{Dm}}}\]

\[\fitmath{\mathrm{nb}=\frac{1}{(0{,}35)}\approx 2{,}857143\,}\]

\item Fracción volumétrica de fibra.

\[\fitmath{\mathrm{Vf}=0{,}55\approx 0{,}55\,}\]

\item Módulo de la fibra.

\[\fitmath{\mathrm{Ef}=200\approx 200\,\mathrm{GPa}}\]

\item Módulo de la matriz.

\[\fitmath{\mathrm{Em}=4\approx 4\,\mathrm{GPa}}\]

\item Módulo longitudinal por la regla de mezclas.

\[\fitmath{\mathrm{EL}=\mathrm{Vf}\cdot \mathrm{Ef}+\left(1-\mathrm{Vf}\right)\cdot \mathrm{Em}}\]

\[\fitmath{\mathrm{EL}=(0{,}55)\cdot (200)+\left(1-(0{,}55)\right)\cdot (4)\approx 111{,}8\,\mathrm{GPa}}\]

\item Módulo transversal por el modelo de Reuss.

\[\fitmath{\mathrm{ET}=\frac{1}{\frac{\mathrm{Vf}}{\mathrm{Ef}}+\frac{1-\mathrm{Vf}}{\mathrm{Em}}}}\]

\[\fitmath{\mathrm{ET}=\frac{1}{\frac{(0{,}55)}{(200)}+\frac{1-(0{,}55)}{(4)}}\approx 8{,}67679\,\mathrm{GPa}}\]

\item Densidad del compuesto.

\[\fitmath{\rho=\mathrm{Vf}\cdot 1800+\left(1-\mathrm{Vf}\right)\cdot 1200}\]

\[\fitmath{\rho=(0{,}55)\cdot 1800+\left(1-(0{,}55)\right)\cdot 1200\approx 1530\,\mathrm{kg/m^{3}}}\]

\item Módulo longitudinal específico.

\[\fitmath{\mathrm{ELesp}=\frac{\mathrm{EL}\cdot 1000000000}{\rho}}\]

\[\fitmath{\mathrm{ELesp}=\frac{(111{,}8)\cdot 1000000000}{(1530)}\approx 7{,}30719\times10^{7}\,\mathrm{J/kg}}\]

\item Tenacidad a fractura en MPa por raíz de metro.

\[\fitmath{K1c=40\approx 40\,}\]

\item Factor geométrico de grieta superficial.

\[\fitmath{Y=1{,}12\approx 1{,}12\,}\]

\item Tensión de fallo con grieta de 4 mm.

\[\fitmath{\mathrm{sc}_{1}=\frac{K1c}{Y\cdot \sqrt{\pi\cdot 0{,}004}}}\]

\[\fitmath{\mathrm{sc}_{1}=\frac{(40)}{(1{,}12)\cdot \sqrt{\pi\cdot 0{,}004}}\approx 318{,}5936\,\mathrm{MPa}}\]

\item Tensión de fallo con grieta de 8 mm.

\[\fitmath{\mathrm{sc}_{2}=\frac{K1c}{Y\cdot \sqrt{\pi\cdot 0{,}008}}}\]

\[\fitmath{\mathrm{sc}_{2}=\frac{(40)}{(1{,}12)\cdot \sqrt{\pi\cdot 0{,}008}}\approx 225{,}2797\,\mathrm{MPa}}\]

\item Tamaño crítico a 250 MPa.

\[\fitmath{\mathrm{ac}=\frac{{\left(\frac{K1c}{Y\cdot 250}\right)}^{2}}{\pi}\cdot 1000}\]

\[\fitmath{\mathrm{ac}=\frac{{\left(\frac{(40)}{(1{,}12)\cdot 250}\right)}^{2}}{\pi}\cdot 1000\approx 6{,}49612\,\mathrm{mm}}\]

\end{enumerate}
